\section{总结与延伸}
\begin{table}[H]
\centering
\begin{tabular}{p{4.5cm}p{11cm}}
\toprule
主题 & 讲座要点 \\
\midrule
自我奖赏架构 & Task Generator 生成多样任务，Solution Generator 输出 Chain-of-Thought，Reward Model 通过 verified response 自我评分。 \\
Self-Instruct + Verified Response & Self-Instruct 确保任务宽度，verified response 提供 cross-check 样本，防止 chaining 过度。 \\
Meta-Judge 与评估 & Meta-Judge 用 SelfT evaluators 模拟 judge，能区分 draft 与 verified response，有效突破 plateau。 \\
部署挑战 & inference-time compute 需可控，hallucination 需要 multi-evaluator，self-aware 系统能更合理地寻求人类介入。 \\
\bottomrule
\end{tabular}
\caption{本讲 self-improving pipeline 核心脉络。}
\end{table}

\begin{table}[H]
\centering
\begin{tabular}{p{5cm}p{10cm}}
\toprule
研究方向 & 说明 \\
\midrule
可验证奖励 & 构造 verified response，让 reward 模型在无需人工标注的情况下保持可信度。 \\
SelfT + Meta Judge & 用 SelfT evaluators 训练 meta judge，从 draft 与 verified response 中自动恢复 high-quality judging。 \\
在线交互 & 通过 Self-Feeding 机制把真实交互、verified response 与 instruction data 串联，形成持续 learning loop。 \\
\bottomrule
\end{tabular}
\caption{后续应持续关注的 research road map。}
\end{table}

\subsection{拓展阅读}
\begin{itemize}
    \item OpenAI “Self-Feeding Chatbot” 论文，展示 verified response 在对话中的迭代。
    \item Weston 等的 SelfT Evaluators / Self-Rewarding 研究，详细解读 judge training。
    \item Alpaca Eval 2 benchmark，提供 win rate 标杆，配合 meta judge 验证。
    \item Chain-of-Thought、RLHF/DPO 的对比文章，有助判断何时使用 self-rewarding 。
\end{itemize}

\subsection{本章小结}
Jason Weston 这堂讲座告诉我们，推理能力是一条需要 ``提出问题、验证答案、判断结果'' 三步走的链子。只有让模型自己从 verified response 中学习、在 meta judge 的引导下再次确认，才能把 reasoning 训练带到下一个台阶。

\end{document}
